\subsection{Historical context}
\label{subsec:history}

The nonordering theorem of Mermin and Wagner, and Mermin's classical
version, exclude spontaneous magnetization in the corresponding
finite-range two-dimensional systems~\cite{MW,Mermin1967}.
McBryan and Spencer obtained algebraic upper bounds on spin correlations
by complex rotations~\cite{McBryanSpencer}.  Such an upper bound does not
distinguish algebraic decay from exponential decay.  Local Ward
identities gave further correlation inequalities and high-temperature
mass-gap criteria in the work of Aizenman and
Simon~\cite{AizenmanSimonWard}.  A continuum mass theorem additionally
requires control as the bare inverse coupling diverges and the lattice
spacing vanishes.

For more than two spin components, perturbative renormalization
predicts that weak microscopic coupling produces an exponentially
large correlation length.  Polyakov identified the renormalization-group
mechanism for two-dimensional Goldstone fields~\cite{Polyakov}.
Br\'ezin and Zinn-Justin developed the expansion near two dimensions,
and Br\'ezin, Zinn-Justin and Le Guillou analyzed the renormalization of
the model and its composite operators~\cite{BrezinZinnJustin,BrezinZinnJustinLeGuillou}.
These works explain why weak microscopic coupling can coexist with a
finite long-distance mass scale.  They also identify the perturbative
coefficients against which a constructive renormalization should be
checked.  Establishing those coefficients is logically different from
bounding the nonperturbative remainder throughout a trajectory whose
length diverges with the cutoff.

Integrability gives a complementary and much more detailed picture of
the expected continuum theory.  Zamolodchikov and Zamolodchikov solved
the factorization, crossing and unitarity equations for
$O(N)$-invariant scattering and identified the sigma-model
solution~\cite{ZamolodchikovZamolodchikov}.  Hasenfratz, Maggiore and
Niedermayer obtained the exact mass-to-renormalization-scale relation
for $O(3)$ and $O(4)$ by matching the Bethe ansatz with perturbation
theory~\cite{HasenfratzMaggioreNiedermayer}; Hasenfratz and Niedermayer
extended the calculation to arbitrary $N\ge3$~\cite{HasenfratzNiedermayer}.
Those calculations determine a mass ratio within the integrable
continuum description.  The additional problem addressed here is to
construct the field distributions from the specified nearest-neighbor
measure and to establish their spectral properties after removal of
the cutoff.  No factorized scattering matrix is used as a premise of
the construction.

Susceptibility and second-moment correlation length provide intrinsic
field and length units.  This is the familiar zero-momentum normalization,
used explicitly by Campostrini, Pelissetto, Rossi and
Vicari~\cite{CampostriniPelissettoRossiVicari1997}.
Their strong-coupling calculations and the lattice--bootstrap comparisons
of Balog and collaborators~\cite{BalogEtAl1999} studied normalized
amplitudes and supplied evidence for continuum universality.
Here the question is convergence of the entire canonically normalized
Schwinger hierarchy for the specified nearest-neighbor model and periodic
thermodynamic state, as the bare coupling tends to infinity without
restriction to a selected sequence.  Comparing this limit with a
terminal-coupling construction requires control of both the field
distributions and the normalizing correlation moments.

Rigorous work has supplied several distinct parts of this program.
Kupiainen proved asymptotic validity of the $1/N$ expansion for lattice
models above the spherical-model critical temperature and obtained a
mass gap for sufficiently large $N$ at each such
temperature~\cite{Kupiainen}.
Kopper later proved mass generation at sufficiently large finite $N$
with a suitable ultraviolet cutoff~\cite{Kopper}.  These large-$N$
results do not specialize to a cutoff-removed construction at $N=3$.
Gaw\k{e}dzki and Kupiainen constructed a continuum limit with asymptotic
freedom for a sigma model with a \emph{hierarchical} kinetic
term~\cite{GawedzkiKupiainen}.  In that setting the hierarchical covariance reduces the flow to
a single-spin recursion.  For the nearest-neighbor model, integration
generates interactions between blocks whose spatial decay must be
controlled.
Mitter and Ramadas constructed the Wilson renormalized trajectory in
perturbation theory in the effective charge and developed the
corresponding effective-action analysis~\cite{MitterRamadas}.

The organization by successive block integrations belongs to the
renormalization-group tradition of Kadanoff and
Wilson~\cite{Kadanoff,Wilson}.  Ba\l aban developed localized effective actions and constrained
variational methods in lattice gauge theory~\cite{BalabanGaugeFlow,BalabanGaugeVariational}
and in the low-temperature analysis of classical $N$-vector
models~\cite{BalabanFlow,BalabanVariational}.
Goswami studied the constrained background-field problem for
two-dimensional $SU(n)$ chiral models, including the $O(4)$
case~\cite{Goswami2024}.
For the two-dimensional $O(4)$ model, Dybalski, Stottmeister and
Tanimoto proved existence and uniqueness for the constrained
small-field variational problem associated with one block-averaging
operation~\cite{DybalskiStottmeisterTanimotoVariational}.
Aru, Garban and Sep\'ulveda proved that, after fixing one spin
direction, rescaled transverse fluctuations converge locally as the
inverse coupling tends to infinity to a rooted vector-valued lattice
Gaussian free field~\cite[Theorem~1.3]{AGS}.  For $O(3)$ this field has
two components.  Their result identifies the local Gaussian
approximation; the construction below also requires control at
spatial scales that grow with the inverse coupling.
